\subsection{The choice between embedded cuts and geodesic cuts}
\label{sec:A1-cut-choice}

\begin{proposition}[Refuted: arbitrary analytic cuts canonically give toric cap identifications]
\label{A1-nontoric-cut}
If the cut circles in 67/074 are interpreted as arbitrary embedded
smooth, or even analytic, circles, cutting and canonical capping
does not always give the asserted antiisomorphisms of pointed caps.
This failure already occurs on a four-punctured sphere with one
separating circle. The geodesic-cut model used in
Theorems~\ref{A1-graph-range} and~\ref{A1-bordered-shadow-range}
therefore requires either geodesic representatives or an explicit
toric compatibility condition; smoothness alone does not supply it.
\end{proposition}
\begin{proof}
Take two hyperbolic pants, each with two cusp ends and one
geodesic boundary of length $2\pi$. Let $w_1,w_2$ be their canonical
cap coordinates, so that $w_j=e^{i\theta_j}$ parametrizes the
boundary by normalized arclength. Glue these boundaries by
\[
 w_2=h(w_1),\qquad
 h(e^{i\theta})=
       \exp\bigl(i(-\theta+\tfrac12\sin\theta)\bigr).
\]
The angular derivative is $-1+\tfrac12\cos\theta$, everywhere
between $-3/2$ and $-1/2$. Thus $h$ is an orientation-reversing
analytic diffeomorphism of the circle. It extends to the holomorphic
annulus map
\[
 H(w)=w^{-1}\exp\bigl((w-w^{-1})/4\bigr).
\]
Its derivative does not vanish on the circle. Since its restriction
to that circle is injective, compactness and local injectivity imply
that it is injective on a sufficiently small whole annulus. Gluing by
this collar extension gives a Riemann surface with an analytic
seam, topologically a four-punctured sphere. Equivalently this
is the analytic special case of Theorem~\ref{A3-welding}.

Cutting along that seam recovers the two original conformal
pants. Their unique canonical hyperbolic metrics, and hence their
canonical cap coordinates, are exactly those used above. An
antiisomorphism of centered unit discs has the form
$w\mapsto c\overline w$, $|c|=1$, by the Schwarz lemma. Its boundary
angular derivative is constantly $-1$. The identification $h$
has derivatives $-1/2$ at $\theta=0$ and $-3/2$ at $\theta=\pi$,
so it is not the boundary map of any such pointed-disc
antiisomorphism. The tensor claimed in 67/075 cannot therefore
be extracted by that stated construction from this analytic cut.
\end{proof}

If the cut datum means an essential isotopy class instead, the
canonical hyperbolic metric of the entire surface selects its
unique geodesic representative. Simultaneously disjoint classes
have disjoint geodesic representatives. Cutting along those
representatives restricts the same hyperbolic metric to both
banks, so the gluing is an arclength isometry. This is the
geodesic interpretation for which the exact range theorems
above apply. It changes the embedded analytic seam of the
counterexample and must be named as part of the definition.
